\documentclass[10pt,letterpaper]{article}
\usepackage{technical-testing}
\setDocTitle{Frame Time And Exposure Problem Set}
\setDocSubtitle{Original problems with numerical checks and evidence requirements}
\setDocVersion{1.0}
\setDocStatus{Educational / Draft}
\setDocOwner{Game Development}
\setDocDate{2026-10-30}
\setDocSummary{Original problems with numerical checks and evidence requirements}
\begin{document}
\makedoctitle
\makedocsummary

\section{Worked conventions}
Display-frame indices are zero based. Intervals are half open; output sampling uses $t_k=k/F$. A complete drawing exposure covers one or more consecutive display indices.
\section{Problems}
\begin{enumerate}
\item A 96-frame scene at 24 fps has a 12-frame anticipation, 30-frame main motion, and the rest as settle and hold. Calculate all boundary indices and durations.
\item A 150-frame sequence at $30000/1001$ fps has a dialogue cue at 2.0 s. Identify the first display frame at or after the cue, stating the rounding policy.
\item Convert 48 display frames of a twos assignment to unique cel exposures. Replace six display frames with ones while preserving the total length. Count the new distinct cels.
\item Design a failure test for an exposure-sheet parser that has a two-frame gap.
\end{enumerate}
\section{Solution checks}
For Problem 1, total duration is 4 seconds, with boundaries 0, 12, 42 and 96. For Problem 2, $2\cdot30000/1001\approx59.94$, so the first frame at or after 2.0 s is frame 60. For Problem 3, original is 24 drawings, replacing three pairs by six single exposures gives 27 drawings if no drawings are intentionally reused.
\section{Assessment}
Give the formula, numerical result, exact time base, a visual timeline, and at least one edge case for every computation.
\end{document}
